\documentclass[twocolumn,a4paper,11pt]{article}

\usepackage[margin=1.8cm]{geometry}
\usepackage{amsmath,amssymb,amsthm}
\usepackage{mathtools}
\usepackage{microtype}
\usepackage{hyperref}
\usepackage{cleveref}
\usepackage{enumitem}
\usepackage{booktabs}
\usepackage{array}
\usepackage{xcolor}
\usepackage[numbers,sort&compress]{natbib}

% Theorem environments
\theoremstyle{plain}
\newtheorem{theorem}{Theorem}[section]
\newtheorem{lemma}[theorem]{Lemma}
\newtheorem{corollary}[theorem]{Corollary}
\newtheorem{proposition}[theorem]{Proposition}

\theoremstyle{definition}
\newtheorem{definition}[theorem]{Definition}
\newtheorem{axiom}[theorem]{Axiom}

\theoremstyle{remark}
\newtheorem{remark}[theorem]{Remark}
\newtheorem{example}[theorem]{Example}

% Math operators
\DeclareMathOperator{\hol}{hol}
\DeclareMathOperator{\diam}{diam}
\DeclareMathOperator{\OR}{OR}
\DeclareMathOperator{\val}{val}
\DeclarePairedDelimiter{\abs}{\lvert}{\rvert}
\DeclarePairedDelimiter{\norm}{\lVert}{\rVert}

\newcommand{\Val}{\mathrm{Val}}
\newcommand{\Sflat}{\mathcal{S}_{\flat}}

\title{\textbf{Wind Tunnel: A Global Self-Consistency Framework\\
for Software Correctness}}

\author{Anonymous}
\date{}

\begin{document}

\maketitle

\begin{abstract}
Contemporary software testing identifies local failures --- wrong syntax,
wrong runtime output --- but is structurally blind to global semantic
failure: a system in which every component passes every test yet fails its
purpose. We formalise this \emph{local invisibility problem} and prove it
is not a limitation of engineering practice but a theorem about bounded
observers. We introduce the \emph{wind tunnel} framework, which replaces
the binary pass/fail oracle with a scalar \emph{semantic entropy}
$\mathcal{S}$, a positive-volume \emph{action-cell} of purpose, and an
ensemble \emph{synchronisation order parameter} $R_\mathrm{ens}$ derived
from Kuramoto coupling theory. We prove three foundational results: the
\emph{Local Invisibility Theorem} (local validity does not imply global
correctness), the \emph{No Template Theorem} (no bounded observer can
encode a complete correctness criterion for a system of equal or greater
complexity), and the \emph{Cycle Inconsistency Theorem} (nonzero holonomy
on any directed cycle implies global semantic incorrectness; zero holonomy
is necessary but not sufficient for global correctness).
We derive a five-regime classification of coordination states and a
two-phase testing protocol whose primary output is a regime map over the
ensemble dependency graph. The framework subsumes unit testing,
property-based testing, and chaos engineering as special cases and provides
the first principled criterion for \emph{purposelessness detection}.
\end{abstract}

%===========================================================================
\section{Introduction}
\label{sec:intro}
%===========================================================================

Dijkstra observed in 1972 that ``testing shows the presence, not the
absence, of bugs'' \cite{dijkstra1972}. The remark is universally quoted
but insufficiently generalised. It is standardly interpreted as an
engineering caveat: test suites are finite, so they miss cases. The deeper
reading is structural: a test is a \emph{local query} addressed to a
bounded subsystem. It cannot, in principle, verify the global circuit.

Consider the following scenario. An enterprise system comprising twelve
services passes all 4\,200 unit tests, all 380 integration tests, and all
property-based tests generated by a QuickCheck-style harness
\cite{claessen2000}. It then fails in production because the composition
of individually correct data-transformation steps accumulates a bias that
inverts a financial decision rule under load. No single service
misbehaved. The failure resided in the cycle, not in any node.

This failure mode is not an edge case. It is the generic failure mode for
systems of sufficient complexity, because generic systems have global
constraints that no local query can reach. The present paper makes this
claim precise.

\subsection{The Wind Tunnel Analogy}

An aeronautical wind tunnel does not test individual rivets. It tests the
\emph{assembled aircraft} at operational load, measuring lift, drag, and
turbulence --- quantities that are properties of the whole, not of any
part. There is no ``expected output'' written on a card. The criterion is
self-consistency under load: does the airflow cohere with the structure?
The tunnel reveals global failure modes invisible to component inspection.

We argue that software correctness requires an analogous instrument. The
wind tunnel framework developed here replaces the binary oracle with:
\begin{enumerate}[leftmargin=*]
  \item A scalar \emph{semantic entropy} $\mathcal{S} \in [0, \Sigma]$
        measuring residual distance from purposeful operation.
  \item A positive-volume \emph{action-cell} $C^* \subseteq X$ in outcome
        space --- the region, not the point, that constitutes correct
        behaviour.
  \item An ensemble synchronisation parameter $R_\mathrm{ens} \in [0,1]$
        measuring phase coherence across the dependency graph.
\end{enumerate}
The primary output of a wind tunnel run is not a verdict but a
\emph{regime map}: a classification of every sub-ensemble by its
coordination state, with decoherence zones identified.

\subsection{Contributions}

The paper makes the following formal contributions:
\begin{itemize}[leftmargin=*]
  \item Proof of the \emph{Local Invisibility Theorem}
        (\cref{thm:local-invisibility}).
  \item Proof of the \emph{No Template Theorem}
        (\cref{thm:no-template}).
  \item Definition of holonomy-based self-consistency and proof of the
        \emph{Cycle Inconsistency Theorem} (\cref{thm:hac}).
  \item Formal definition of semantic entropy $\mathcal{S}$ and the
        action-cell $C^*$, with floor-positivity and cell-truth
        established as principled axioms.
  \item Kuramoto-derived five-regime classification and the
        \emph{Synchronisation as Partition Extinction} theorem
        (\cref{thm:partition-extinction}).
  \item The two-phase wind tunnel testing protocol and a formal criterion
        for purposelessness detection.
\end{itemize}

\subsection{Relation to Existing Work}

Unit testing \cite{myers1979} and model checking \cite{clarke1981} operate
via local queries and template comparison respectively; both are shown to
be special limiting cases of wind tunnel testing. Property-based testing
\cite{claessen2000} generalises the test-case lattice but not the observer
bound. Chaos engineering \cite{basiri2016} injects perturbations into live
systems but lacks a formal correctness criterion. Formal verification
\cite{clarke1981} provides completeness within a finite specification, but
the specification is itself a bounded observer (\cref{sec:no-template}).
Wind tunnel testing is the first framework to ground correctness in
self-consistency rather than template matching.

%===========================================================================
\section{Foundations}
\label{sec:foundations}
%===========================================================================

We establish the formal vocabulary from first principles.

\subsection{Outcome Space and Trajectories}

\begin{definition}[Outcome Space]
\label{def:outcome-space}
Let $X$ be a metric space $(X, d_X)$ of \emph{observable outcomes}. Each
point $x \in X$ represents a complete, fully realised state of a running
system at a given instant: all variable values, all message buffers, all
resource allocations.
\end{definition}

The outcome space is not the state space of any single unit; it is the
product space of all unit states under all possible external inputs. Even
for programs operating over finite alphabets, this space is in general
uncountably infinite because possible execution interleavings grow
super-exponentially with the number of units.

\begin{definition}[Trajectory and Terminus]
\label{def:trajectory}
A \emph{trajectory} is a continuous map $\gamma: [0,T] \to X$ describing
system evolution from initial state $\gamma(0)$ over time horizon $T$.
The \emph{terminus} is the terminal point $\Gamma_f = \gamma(T)$.
\end{definition}

\begin{definition}[Computational Unit]
\label{def:unit}
A \emph{computational unit} $u$ is a tuple $(A_u, P_u, \mathrm{val}_u)$
where:
\begin{itemize}[leftmargin=*]
  \item $A_u \subseteq X$ is the \emph{aperture}: the set of outcomes $u$
        can process (its categorical input domain).
  \item $P_u \subseteq X$ is the \emph{production set}: outcomes $u$ can
        produce (its output range).
  \item $\mathrm{val}_u: X \to \{0,1\}$ is the \emph{local validation
        function}: $\mathrm{val}_u(x) = 1$ iff $u$'s contribution to
        state $x$ is locally consistent.
\end{itemize}
\end{definition}

\begin{remark}
The distinction between $A_u$ and $P_u$ is the distinction between what
a unit can receive and what it can emit. An aperture mismatch ---
$P_{u_i} \not\subseteq A_{u_j}$ for adjacent units $u_i, u_j$ --- is a
categorical interface failure detectable statically. Aperture compatibility
is necessary but not sufficient for global correctness.
\end{remark}

\subsection{The System Graph}

\begin{definition}[System Graph]
\label{def:system-graph}
A \emph{system} $S$ is a directed graph $G = (V, E)$ where $V = \{u_1,
\ldots, u_n\}$ is a finite set of computational units and
$(u_i, u_j) \in E$ iff $P_{u_i} \cap A_{u_j} \neq \emptyset$ (unit
$u_i$ produces outputs that unit $u_j$ can process).
\end{definition}

\begin{definition}[Global Validation]
\label{def:global-val}
The \emph{global validation} of system $S$ at state $x \in X$ is
$\Val_S(x) = \prod_{u \in V} \mathrm{val}_u(x)$. We say $S$ \emph{passes
local validation} at $x$ if $\Val_S(x) = 1$.
\end{definition}

\subsection{Semantic Entropy}

\begin{definition}[Semantic Entropy]
\label{def:s-entropy}
The \emph{semantic entropy functional} $\mathcal{S}: X \to [0, \Sigma]$
assigns to each outcome $x \in X$ a non-negative real number measuring
the residual semantic distance of $x$ from purposeful operation. The
constant $\Sigma > 0$ is the \emph{maximal semantic distance} (canonical
normalisation $\Sigma = 100$).
\end{definition}

\begin{axiom}[Monotonicity]
\label{ax:monotone}
If $x$ and $x'$ differ only in that $x'$ is measurably closer to purpose
by any observable criterion, then $\mathcal{S}(x') \leq \mathcal{S}(x)$.
\end{axiom}

\begin{axiom}[Continuity]
\label{ax:continuity}
$\mathcal{S}$ is Lipschitz continuous: there exists $L > 0$ such that
$\abs{\mathcal{S}(x) - \mathcal{S}(x')} \leq L \cdot d_X(x, x')$ for
all $x, x' \in X$.
\end{axiom}

\begin{axiom}[Floor Positivity]
\label{ax:floor}
For every bounded computational system $S$, the infimum
$\Sflat(S) = \inf_{x \in X} \mathcal{S}(x) > 0$.
\end{axiom}

\begin{remark}[Motivation for Floor Positivity]
Axiom~\ref{ax:floor} reflects the operational irreducibility of finite
observers: any bounded system evaluating its own state incurs a nonzero
residual from self-reference and finite representational capacity. We take
this as an axiom; it is not derived from G\"{o}del or Shannon but is
independently motivated by the observer-boundedness of any real measurement
system. A bounded computational system cannot measure its own state with
zero uncertainty, and the floor captures precisely this irreducible
residual. The axiom is consistent with, but not entailed by, those deeper
results.
\end{remark}

\subsection{The Action-Cell}

\begin{definition}[Action-Cell]
\label{def:action-cell}
The \emph{action-cell} $C^* \subseteq X$ of a system with purpose $\phi$
is the maximal open set satisfying:
\begin{enumerate}[leftmargin=*]
  \item $\mathcal{S}(x) = \Sflat$ for all $x \in C^*$ (the cell achieves
        the floor).
  \item $C^*$ is path-connected in $X$.
  \item $\diam(C^*) > 0$ (the cell has positive volume).
\end{enumerate}
\end{definition}

\begin{remark}[Cell-Truth]
The cell-truth principle holds: correctness is not a point but a region of
positive volume. Two outcomes $x, x'$ with $d_X(x,x') > 0$ can both lie
in $C^*$ and be semantically indistinguishable, because $\mathcal{S}$ is
flat on $C^*$ by construction. This directly contradicts the assumption
underlying point-oracle testing, which requires $C^*$ to be a singleton.
The cell has positive measure precisely because perfect knowledge is
forbidden by the floor positivity axiom: the floor value $\Sflat > 0$
means the cell is optimal relative to the ensemble's bounded capacity,
not optimal in any absolute sense.
\end{remark}

%===========================================================================
\section{The Local Invisibility Theorem}
\label{sec:local-invisibility}
%===========================================================================

\begin{definition}[Global Semantic Correctness]
\label{def:global-correctness}
A system $S$ is \emph{globally semantically correct} at state $x$ if
$x \in C^*$, i.e., $\mathcal{S}(x) = \Sflat$.
\end{definition}

\begin{theorem}[Local Invisibility Theorem]
\label{thm:local-invisibility}
There exist systems $S = (G, \{\mathrm{val}_u\}_{u \in V})$ and states
$x \in X$ such that $\Val_S(x) = 1$ yet $x \notin C^*$.
\end{theorem}

\begin{proof}
We construct an explicit witness. Let $V = \{u_1, u_2, u_3\}$ with
$E = \{(u_1,u_2),(u_2,u_3),(u_3,u_1)\}$, forming a directed cycle
$c = (u_1, u_2, u_3, u_1)$.

Let the local state of each unit be a real scalar $x_i \in \mathbb{R}$,
so the global state is $x = (x_1, x_2, x_3) \in \mathbb{R}^3$.
Define:
\[
  \mathrm{val}_{u_i}(x) = \mathbf{1}[x_i \in [0,1]], \quad i \in
  \{1,2,3\}.
\]

Let each edge $(u_i, u_{i+1 \bmod 3})$ apply a transformation
$f: t \mapsto t + \delta$ for a fixed $\delta > 0$. After traversing
the full cycle once starting from $x_1$, the output returned to $u_1$
is $x_1 + 3\delta$. For the loop to close, we would need $3\delta = 0$,
i.e., $\delta = 0$.

Set $\delta = \varepsilon > 0$ and $x_1 = x_2 = x_3 = 0.3$. Then
$x_i \in [0,1]$ for all $i$, so $\mathrm{val}_{u_i}(x) = 1$ for all
$i$ and thus $\Val_S(x) = 1$.

Yet the cycle accumulates error $3\varepsilon$ per traversal. After $k$
traversals: $x_1^{(k)} = 0.3 + 3k\varepsilon$. For
$k \geq \lceil 0.7 / (3\varepsilon) \rceil$, the state exits $[0,1]$
--- a condition invisible at the initial snapshot but structurally
guaranteed.

To see that $x \notin C^*$: the purpose cell $C^*$ requires holonomy-free
closure of every constraint cycle (formalised as \cref{def:holonomy}). The
state $x$ with $\delta = \varepsilon > 0$ has nonzero holonomy
$\hol(c, x) = 3\varepsilon > 0$. By the Cycle Inconsistency Theorem
(\cref{thm:hac}), nonzero holonomy implies $x \notin C^*$. Therefore
$\Val_S(x) = 1$ and $x \notin C^*$ simultaneously. \qed
\end{proof}

\begin{remark}[Circuit Analogy]
The proof is directly analogous to Kirchhoff's voltage law
\cite{kirchhoff1847}. In an electrical network, every component may read
correct voltage locally while the directed sum of voltage drops around a
closed loop is nonzero --- a violation detectable only by closing the loop.
The local validation function $\mathrm{val}_{u_i}$ is a ``local voltmeter''
reading correctly at each node, while the cycle residual $3\varepsilon$ is
a holonomy that no finite set of local measurements can detect.
\end{remark}

\begin{corollary}
\label{cor:unit-test-blind}
Any test suite composed exclusively of per-unit queries is structurally
blind to the failure class exhibited in \cref{thm:local-invisibility}.
\end{corollary}

\begin{proof}
A unit test for unit $u_i$ evaluates $\mathrm{val}_{u_i}(x)$ with all
other units mocked or absent. The accumulation residual $\delta$ in the
constructed witness is a property of the cycle $E$: it does not appear in
any unit's local state at a single snapshot. Hence no unit test, regardless
of how many inputs it samples, can observe the inconsistency. \qed
\end{proof}

%===========================================================================
\section{The No Template Theorem}
\label{sec:no-template}
%===========================================================================

Unit testing assumes an oracle --- a reference against which output is
compared. The No Template Theorem shows this assumption fails beyond a
complexity threshold.

\begin{definition}[Bounded Observer]
\label{def:bounded-observer}
A \emph{bounded observer} $R$ is a system with finite state capacity
$\abs{K_R} < \infty$: the number of distinct states representable by
$R$ is bounded.
\end{definition}

\begin{definition}[State Complexity]
\label{def:state-complexity}
The \emph{state complexity} $\kappa(S)$ of a system $S$ is
$\kappa(S) = \log_2 \abs{\{\text{reachable states of } S\}}$.
\end{definition}

\begin{theorem}[No Template Theorem]
\label{thm:no-template}
Let $R$ be a bounded observer with state capacity $\abs{K_R}$, and let
$S$ be a system with $\kappa(S) > \log_2 \abs{K_R}$. Then there is no
computable function $\phi_R: X \to \{0,1\}$ stored in $R$ such that
$\phi_R(x) = 1$ if and only if $x \in C^*$.
\end{theorem}

\begin{proof}
A function $\phi_R: X \to \{0,1\}$ correctly identifying $C^*$ must
distinguish $x \in C^*$ from $x' \notin C^*$ for every $x, x' \in X$.
Since $C^*$ is a subset of $S$'s reachable state space, the number of
binary decisions $\phi_R$ must encode is at least $2^{\kappa(S)}$.

However, $\phi_R$ is computed by $R$, which has state capacity $\abs{K_R}$.
By the pigeonhole principle, $R$ can represent at most $\abs{K_R}$ distinct
decision procedures. Computing $\phi_R$ requires at least one internal
state of $R$ to distinguish each pair $(x, x')$ with $x \in C^*$ and
$x' \notin C^*$.

When $\kappa(S) > \log_2 \abs{K_R}$, the number of such pairs exceeds
$\abs{K_R}$. Therefore no function stored in $R$ can correctly decide
membership in $C^*$ for all reachable states of $S$. \qed
\end{proof}

\begin{corollary}
\label{cor:test-suite-incomplete}
Every test suite of bounded complexity is an incomplete specification of
correctness for any system of equal or greater complexity.
\end{corollary}

\begin{proof}
A test suite is a bounded observer: it has finitely many test cases, each
encoding a finite input-output pair, giving finite state capacity. If the
system under test has state complexity exceeding the test suite's
$\log_2 \abs{K_R}$, the result follows directly from \cref{thm:no-template}.
\qed
\end{proof}

\begin{remark}
Corollary~\ref{cor:test-suite-incomplete} is qualitatively distinct from
the classical observation that test suites cannot cover all inputs
\cite{dijkstra1972}. That observation is quantitative: more tests reduce
uncovered cases. Our result is representational: no increase in test count
can encode a complete correctness criterion when the system's state
complexity exceeds the observer's capacity. Fisher's theory of sufficient
statistics \cite{fisher1925} provides the information-geometric
interpretation: a test suite is a sufficient statistic for correctness only
when its Fisher information matrix spans the system's state manifold ---
a condition structurally violated when $\kappa(S) > \log_2 \abs{K_R}$.
\end{remark}

%===========================================================================
\section{Self-Consistency as Correctness}
\label{sec:hac}
%===========================================================================

Since template matching is impossible beyond the complexity threshold of
\cref{sec:no-template}, we require an alternative definition of
correctness. We provide it via holonomy.

\subsection{Holonomy in the Dependency Graph}

\begin{definition}[Constraint Path]
\label{def:constraint-path}
A \emph{constraint path} $\pi = (e_1, \ldots, e_k)$ is a sequence of
edges in $G$ forming a walk. The \emph{composed transformation}
$T_\pi: X \to X$ is the sequential composition of unit transformations
along $\pi$.
\end{definition}

\begin{definition}[Cycle Specification]
\label{def:cycle-spec}
Each directed cycle $c \in \mathcal{C}$ carries a \emph{specification}
$T_c^{\mathrm{spec}}: X \to X$ --- the transformation the system designer
intends the cycle to perform. This specification is part of the system's
purpose functional. For stateless, idempotent cycles, $T_c^{\mathrm{spec}}
= \mathrm{id}$. For stateful cycles (accumulators, iterative solvers,
retry loops, state machines), $T_c^{\mathrm{spec}}$ encodes the intended
state evolution.
\end{definition}

\begin{definition}[Holonomy]
\label{def:holonomy}
Let $\mathcal{C}$ be the set of all directed cycles in $G$. The
\emph{holonomy} of cycle $c \in \mathcal{C}$ at state $x$ is:
\[
  \hol(c, x) = \norm{T_c(x) - T_c^{\mathrm{spec}}(x)}_{X},
\]
where $T_c: X \to X$ is the actual composition of transformations around
$c$. Zero holonomy means the cycle does exactly what it was designed to
do. Nonzero holonomy means the cycle deviates from its specification.
System $S$ is \emph{holonomy-free} if $\hol(c, x) = 0$ for all
$c \in \mathcal{C}$ and all reachable $x$.
\end{definition}

\begin{remark}
For purely stateless, idempotent cycles, $T_c^{\mathrm{spec}} = \mathrm{id}$
and holonomy reduces to $\norm{T_c(x) - x}_X$, recovering the familiar
form that checks whether the cycle returns to its starting state. For
stateful cycles --- accumulators, PID controllers, iterative solvers ---
$T_c^{\mathrm{spec}}$ captures the intended evolution, and holonomy detects
deviation from design rather than deviation from rest. This distinction
matters: a legitimate accumulator has $T_c(x) \neq x$ by design; the
original form would incorrectly condemn it as inconsistent.
\end{remark}

\begin{remark}
Holonomy is the software analogue of Kirchhoff's second law
\cite{kirchhoff1847}. In a directed electrical network, the sum of
voltage drops around any closed loop must be zero (conservation of
energy). In a dependency graph, the actual composition of transformations
around any closed path must match the intended composition (conservation
of semantic content relative to design). Nonzero holonomy means the cycle
deviates from its specification --- a sign of global inconsistency that no
local measurement can detect.
\end{remark}

\subsection{Cycle Inconsistency Detection}

\begin{theorem}[Cycle Inconsistency Theorem]
\label{thm:hac}
If system $S$ has nonzero holonomy on any directed cycle, then $S$ is
globally semantically incorrect: there exist reachable states $x \notin C^*$.
Equivalently, nonzero holonomy on any cycle $c$ implies a detectable
cycle-inconsistency failure. Zero holonomy is necessary but not sufficient
for global correctness.
\end{theorem}

\begin{proof}
We prove the forward direction: nonzero holonomy implies global incorrectness.

Suppose $\hol(c, x) = \delta > 0$ for some directed cycle $c$ and some
reachable state $x$. By definition,
$\norm{T_c(x) - T_c^{\mathrm{spec}}(x)}_X = \delta > 0$, so the actual
cycle transformation deviates from its specification by $\delta$ at $x$.

After one traversal of $c$, the system reaches state $T_c(x)$, which
differs from the intended state $T_c^{\mathrm{spec}}(x)$ by $\delta$.
This deviation is observable: it is a concrete discrepancy between the
system's actual behaviour and its declared specification on this cycle.
A state exhibiting such a discrepancy cannot lie in $C^*$, since $C^*$
is the region where $\mathcal{S}(x) = \Sflat$ --- the floor achievable
only when all constraint relationships are satisfied to within the
representational capacity of the system. A cycle that deviates from its
specification by $\delta > 0$ contributes a positive semantic distance
above $\Sflat$.

Formally: repeated traversals produce $T_c^k(x)$, with accumulated
deviation growing by at least $\delta$ per traversal (the deviation
compounds unless the specification is satisfied). Since $C^*$ has
finite diameter $\diam(C^*)$, the trajectory eventually exits $C^*$,
confirming $x \notin C^*$ under repeated cyclic operation.

The necessity of zero holonomy for global correctness follows immediately:
if any cycle has nonzero holonomy, the above argument produces reachable
states outside $C^*$, so the system is not globally correct.

The sufficiency direction does not hold in general. Zero holonomy on all
cycles means every cycle behaves as specified, but global correctness
additionally requires that the initial state $\gamma(0) \in C^*$ and that
all acyclic paths (boundary conditions) also respect the purpose cell.
DAG-structured systems illustrate this directly: they are trivially
holonomy-free (no directed cycles exist), yet they can be globally incorrect
if their boundary conditions are wrong. Zero holonomy is therefore necessary
but not sufficient for global correctness. \qed
\end{proof}

\begin{remark}
The Cycle Inconsistency Theorem is significant because it provides a
correctness \emph{witness} requiring no external reference for the cycle
behaviour: one need not know the global purpose of the system to detect
cycle-level inconsistency. It suffices to verify that each cycle's actual
transformation matches its declared specification $T_c^{\mathrm{spec}}$.
This mirrors physical law: Kirchhoff's laws assert consistency conditions
without reference to any external ``correct'' state.
\end{remark}

\begin{corollary}[DAG Boundary Condition]
\label{cor:acyclic-invisible}
For directed acyclic graphs, holonomy is trivially zero (no directed cycles
exist), so the Cycle Inconsistency Theorem provides no information about
global correctness. All global correctness verification for DAG-structured
systems must come from \emph{boundary conditions}: the initial states and
terminal outputs. This is a positive result: it tells us precisely where
to look. Holonomy-based methods detect cycle-level failures; boundary
condition methods detect DAG-level failures; together they cover the full
dependency graph.
\end{corollary}

%===========================================================================
\section{Purpose and Its Detection}
\label{sec:purpose}
%===========================================================================

\subsection{Ensembles and Composite Floor}

\begin{definition}[Ensemble]
\label{def:ensemble}
An \emph{ensemble} $E = (G_E, \Phi_E)$ is a system graph $G_E = (V_E, E_E)$
equipped with a \emph{purpose functional}
$\Phi_E: 2^X \to [0, \Sigma]$ assigning semantic entropy to each
candidate action-cell $C \subseteq X$.
\end{definition}

\begin{definition}[Composite Floor]
\label{def:composite-floor}
The \emph{composite floor} of ensemble $E$ is:
\[
  \Sflat(E) = \inf_{\substack{C \subseteq X \\ C \text{ satisfies
  Def.~\ref{def:action-cell}}}} \Phi_E(C).
\]
\end{definition}

\begin{definition}[Purpose Cell]
\label{def:purpose-cell}
A cell $C^* \subseteq X$ is a \emph{purpose cell} for $E$ if
$\Phi_E(C^*) = \Sflat(E)$.
\end{definition}

\subsection{Existence of Purpose}

\begin{theorem}[Purpose Existence Theorem]
\label{thm:purpose-existence}
For every ensemble $E$ with $\Sflat(E) < \Sigma$, there exists a purpose
cell $C^* \subseteq X$.
\end{theorem}

\begin{proof}
Since $\Sflat(E) < \Sigma$ and $\Phi_E$ is continuous (inherited from
the Lipschitz continuity of $\mathcal{S}$, Axiom~\ref{ax:continuity}),
the infimising sequence $\{C_n\}_{n \geq 1}$ with
$\Phi_E(C_n) \to \Sflat(E)$ is non-empty. Equip the space of non-empty
closed subsets of $X$ with the Hausdorff metric $d_H$. By the Banach
fixed-point theorem \cite{banach1922} applied to the complete metric
space $(2^X, d_H)$, every Cauchy sequence converges to a limit point.
The sequence $\{C_n\}$ is Cauchy under the natural ordering induced by
$\Phi_E$, and its limit $C^*$ satisfies $\Phi_E(C^*) = \Sflat(E)$ by
continuity of $\Phi_E$. The path-connectivity and positive-diameter
conditions are preserved in the limit because $\Sflat(E) > 0$
(Axiom~\ref{ax:floor}) prevents degeneration to a singleton. Hence $C^*$
is a valid purpose cell. \qed
\end{proof}

\subsection{Why Isolation Cannot Detect Purposelessness}

\begin{definition}[Purposelessness]
\label{def:purposeless}
Unit $u \in V$ is \emph{purposeless in ensemble $E$} if
$\Sflat(E) = \Sflat(E \setminus \{u\})$: removing $u$ leaves the
composite floor unchanged.
\end{definition}

\begin{theorem}[Isolation Blindness]
\label{thm:isolation-blind}
No function of the form $h: (A_u, P_u, \mathrm{val}_u) \to \{0,1\}$
can correctly decide whether $u$ is purposeless in $E$.
\end{theorem}

\begin{proof}
Purposelessness is a \emph{relational property}. The quantity
$\Sflat(E) - \Sflat(E \setminus \{u\})$ depends on the purpose
functional $\Phi_E$, which integrates contributions of all units through
the full ensemble graph $G_E$.

A function $h$ depending only on $(A_u, P_u, \mathrm{val}_u)$ accesses
only $u$'s local properties. To show this is insufficient, we construct
two ensembles $E_1$ and $E_2$ in which $u$'s local properties are
identical but its purposelessness status differs:

Let $E_1$ be an ensemble where $u$ uniquely provides a function $f$ that
reduces $\Sflat$: $\Sflat(E_1 \setminus \{u\}) > \Sflat(E_1)$, so $u$
is purposeful. Construct $E_2$ from $E_1$ by adding unit $u'$ with
$(A_{u'}, P_{u'}, \mathrm{val}_{u'}) = (A_u, P_u, \mathrm{val}_u)$:
a perfect functional duplicate of $u$ with identical local properties.
In $E_2$, unit $u$ is purposeless:
$\Sflat(E_2 \setminus \{u\}) = \Sflat(E_2)$ because $u'$ provides the
same floor reduction.

Since $(A_u, P_u, \mathrm{val}_u)$ is the same in $E_1$ and $E_2$, any
function $h$ depending only on these properties returns the same value in
both cases, and is therefore wrong in at least one of the two ensembles.
No such $h$ can correctly decide purposelessness. \qed
\end{proof}

\begin{remark}[The Airbus Problem]
\cref{thm:isolation-blind} is the formal statement of a structural
epistemological limit: a worker examined in isolation cannot be assessed
for their contribution to an aircraft's airworthiness. The contribution
exists only through the worker's role in the ensemble. Unit testing commits
this error by design. The theorem shows this is not a practical deficiency
correctable by better tests; it is a structural impossibility.
\end{remark}

%===========================================================================
\section{Coordination Regimes}
\label{sec:kuramoto}
%===========================================================================

The impossibility results of \cref{sec:local-invisibility,sec:no-template,sec:purpose}
establish what testing cannot do in isolation. We now build the positive
theory: a framework for measuring global coordination.

\subsection{Phase Representation and Kuramoto Dynamics}

Drawing on Kuramoto's model of coupled oscillators \cite{kuramoto1975,
kuramoto1984}, we assign each unit $u_i$ an internal \emph{phase}
$\theta_i \in [0, 2\pi)$ representing alignment with the ensemble's
current trajectory.

\begin{definition}[Natural Frequency]
\label{def:natural-freq}
The \emph{natural frequency} $\omega_i$ of unit $u_i$ is the rate at
which $\theta_i$ evolves in isolation, determined by $u_i$'s processing
cycle, data rate, and semantic update rate.
\end{definition}

\begin{definition}[Kuramoto Dynamics]
\label{def:kuramoto}
The \emph{Kuramoto dynamics} of ensemble $E$ on graph $G$ are:
\[
  \dot{\theta}_i = \omega_i + \frac{K}{\deg(u_i)}
  \sum_{(u_i, u_j) \in E} \sin(\theta_j - \theta_i),
  \quad i = 1, \ldots, n,
\]
where $K \geq 0$ is the \emph{coupling strength} and $\deg(u_i)$ is
the in-degree of $u_i$.
\end{definition}

The Kuramoto equation models how units with heterogeneous natural
frequencies are driven toward phase coherence by coupling
\cite{strogatz2000}. When $u_i$'s output feeds $u_j$, $u_i$'s phase
influences $u_j$'s evolution, analogous to coupled oscillators. Vicsek
et al.\ \cite{vicsek1995} demonstrated that analogous local-alignment
rules generate global order above a critical coupling in biological
collective motion systems; the same universality class applies here.

\subsection{The Ensemble Order Parameter}

\begin{definition}[Ensemble Order Parameter]
\label{def:order-param}
The \emph{ensemble order parameter} is:
\[
  R_\mathrm{ens} \cdot e^{i\Psi} = \frac{1}{n} \sum_{k=1}^n
  e^{i\theta_k},
\]
where $R_\mathrm{ens} \in [0,1]$ measures global phase coherence and
$\Psi$ is the mean phase. $R_\mathrm{ens} = 0$ indicates complete
incoherence; $R_\mathrm{ens} = 1$ indicates perfect synchronisation.
\end{definition}

\subsection{Critical Coupling}

\begin{proposition}[Critical Coupling]
\label{prop:critical-coupling}
Let natural frequencies $\{\omega_i\}$ be drawn from a distribution with
standard deviation $\sigma_\omega$. The critical coupling at which
synchronisation onset occurs is:
\[
  K_c = \frac{2\sigma_\omega}{\pi}.
\]
\end{proposition}

\begin{proof}[Proof sketch]
This result is due to Kuramoto \cite{kuramoto1975} and rigorously analysed
by Strogatz \cite{strogatz2000}. The key step linearises the mean-field
self-consistency equation near $R_\mathrm{ens} = 0$. The condition for a
non-trivial solution is $K \int_{-\pi/2}^{\pi/2} g(K r)\cos^2(\theta)\,
d\theta = 1$, where $g$ is the frequency density. For symmetric, unimodal
$g$ (including Gaussian and Lorentzian), this reduces to
$K_c = 2\sigma_\omega / \pi$. \qed
\end{proof}

\begin{remark}
Proposition~\ref{prop:critical-coupling} has a direct engineering
interpretation. High heterogeneity in unit processing speeds (large
$\sigma_\omega$) requires high coupling (tight data contracts, synchronous
protocols) to achieve coherence. Low-coupling architectures (microservices
with loose message passing) can sustain coherence only if units are
sufficiently homogeneous in natural frequency.
\end{remark}

\subsection{Five Coordination Regimes}

\begin{definition}[Coordination Regimes]
\label{def:regimes}
Classify ensemble $E$ by $R_\mathrm{ens}$:
\begin{itemize}[leftmargin=*]
  \item \textbf{Turbulent} ($R_\mathrm{ens} < 0.3$): Semantic incoherence.
        Phase differences dominate. No shared action-cell.
  \item \textbf{Aperture-dominated} ($0.3 \leq R_\mathrm{ens} < 0.5$):
        Units share interface contracts but not purpose. Cluster formation
        without global convergence.
  \item \textbf{Hierarchical cascade} ($0.5 \leq R_\mathrm{ens} < 0.8$):
        Layered coordination. Sub-ensembles cohere internally but not
        globally; decoherent interfaces separate coherent components.
  \item \textbf{Coherent} ($0.8 \leq R_\mathrm{ens} < 0.95$): Ensemble
        converges toward a common action-cell. Semantic entropy declining.
        Residual friction at boundary units.
  \item \textbf{Phase-locked} ($R_\mathrm{ens} \geq 0.95$): Partition
        extinction. All phase differences are constant. Coordination
        friction is zero.
\end{itemize}
\end{definition}

\subsection{Synchronisation as Partition Extinction}

\begin{theorem}[Partition Extinction]
\label{thm:partition-extinction}
The transition from the Coherent to the Phase-locked regime is
discontinuous: coordination friction $\mathcal{F}$ drops to zero in a
single event --- partition extinction --- with no intermediate states of
nonzero-but-small $\mathcal{F}$ in the mean-field limit.
\end{theorem}

\begin{proof}
Define the coordination friction:
\[
  \mathcal{F} = \frac{1}{\abs{E}} \sum_{(u_i,u_j) \in E}
  \bigl(\dot{\theta}_i - \dot{\theta}_j\bigr)^2.
\]

Phase-locking requires every pair $(u_i, u_j) \in E$ to satisfy the
fixed-point condition $K \sin(\theta_j^* - \theta_i^*) = \omega_j -
\omega_i$. This is a system of nonlinear equations. By Landau's theory
of second-order phase transitions \cite{landau1937}, the order parameter
$R_\mathrm{ens}$ near the critical point bifurcates: below $K_c$,
$R_\mathrm{ens} = 0$ is the unique stable fixed point; above $K_c$,
a non-trivial fixed point $R_\mathrm{ens} > 0$ emerges. For the
full locking condition, the analysis follows the BCS theory structure
\cite{bardeen1957}: the system either achieves a globally phase-locked
ground state (zero friction) or it does not.

Formally: a unit $u_i$ is ``locked'' iff
$\dot{\theta}_i = \Omega$ (the common frequency). The set of locked
units $L \subseteq V$ forms the ``condensate.'' As $K$ increases through
$K_c$, the condensate grows. The partition $\{L, V \setminus L\}$ shrinks.
Phase-locking is achieved when $L = V$ --- the partition collapses to
the empty partition. This collapse is instantaneous by the Onsager
variational principle \cite{onsager1953}: the free-energy minimum
transitions discontinuously from $L \subsetneq V$ to $L = V$. There
is no stable intermediate state $L \subsetneq V$ with $\mathcal{F} = 0$,
because such a state would require some units to be locked
($\dot{\theta} = \Omega$) while others are not --- but the locking
condition is global: if any pair $(u_i, u_j) \in E$ is unlocked,
$\mathcal{F} > 0$. \qed
\end{proof}

\begin{corollary}[Subtask Decoupling]
\label{cor:subtask-decoupling}
When $K < K_c$, the local $\mathcal{S}$-value of unit $u$ is completely
decoupled from the global $\mathcal{S}$-value of the ensemble.
\end{corollary}

\begin{proof}
Below $K_c$, each unit evolves near its natural frequency $\omega_i$.
The global $\mathcal{S}$-value is determined by $R_\mathrm{ens} \approx 0$
(Turbulent or Aperture-dominated regime), while each unit's local
validation $\mathrm{val}_{u_i}$ can return 1. A unit can be locally
correct while the ensemble is globally incoherent. Conversely, a unit
with locally poor behaviour is irrelevant to global $\mathcal{S}$ when
$K < K_c$ because its influence does not propagate effectively. Neither
direction of inference is valid. \qed
\end{proof}

%===========================================================================
\section{The Wind Tunnel Protocol}
\label{sec:protocol}
%===========================================================================

We synthesise the preceding results into a concrete testing protocol.

\subsection{Static Phase}

The static phase analyses $G = (V, E)$ and unit specifications without
running the system.

\begin{definition}[Synchronisation Tension]
\label{def:sync-tension}
The \emph{synchronisation tension} between adjacent units $u_i, u_j$
(with $(u_i, u_j) \in E$) is:
\[
  \vartheta(u_i, u_j) = d_X(P_{u_i}, A_{u_j}) +
  \abs{\omega_i^\mathrm{est} - \omega_j^\mathrm{est}},
\]
where $d_X(P_{u_i}, A_{u_j}) = \inf_{p \in P_{u_i}, a \in A_{u_j}}
d_X(p, a)$ is the aperture gap and $\omega_i^\mathrm{est}$ is the
estimated natural frequency from $u_i$'s specification.
\end{definition}

A large aperture gap indicates a categorical interface mismatch: $u_i$
produces outputs $u_j$ cannot process. A large frequency difference
indicates a processing rate mismatch: even if the interface is compatible,
the units operate at semantically incompatible rates.

\begin{definition}[Static Order Parameter Estimate]
\label{def:static-regime}
The \emph{static order parameter estimate} is:
\[
  R_\mathrm{est} = \exp\!\left(-\frac{1}{\abs{E}} \sum_{(u_i,u_j) \in E}
  \vartheta(u_i, u_j)\right).
\]
Classify the ensemble into the regime of Definition~\ref{def:regimes} by
$R_\mathrm{est}$.
\end{definition}

\begin{definition}[Decoherence Zone]
\label{def:decoherence-zone}
A \emph{decoherence zone} is a maximal connected subgraph $G' \subseteq G$
such that $R_\mathrm{est}(G') < R_\mathrm{est}(G)$. Decoherence zones
are sub-ensembles where coordination is locally worse than the global
average.
\end{definition}

\begin{proposition}[Static Phase Outputs]
\label{prop:static-outputs}
The static phase produces: (1)~a regime label for the full ensemble and
each significant sub-ensemble; (2)~a ranked list of decoherence zones by
severity; (3)~a critical coupling estimate $\hat{K}_c = 2\hat\sigma_\omega
/\pi$; and (4)~cycle inconsistency candidates: directed cycles in $G$
whose static accumulated residuals suggest likely holonomy violations.
\end{proposition}

\subsection{Dynamic Phase}

The dynamic phase runs the system under representative load $\Lambda$ and
records trajectories.

\begin{definition}[Runtime Record]
\label{def:runtime-record}
Under load $\Lambda$, the \emph{runtime record} of unit $u$ is
$(\gamma_u, \Gamma_u, M_u)$ where $\gamma_u: [0,T] \to X$ is the
trajectory, $\Gamma_u = \gamma_u(T)$ is the terminus, and $M_u \subseteq X$
is the accumulated memory (set of states visited).
\end{definition}

\begin{definition}[Dynamic Holonomy Measurement]
\label{def:dynamic-holonomy}
For directed cycle $c = (u_{i_1}, \ldots, u_{i_k}, u_{i_1})$, the
\emph{dynamic holonomy} is:
\[
  \hol_\mathrm{dyn}(c) = \norm{\Gamma_{u_{i_k}} -
  T_c^{\mathrm{spec}}(\gamma_{u_{i_1}}(0))}_{X}.
\]
A cycle with $\hol_\mathrm{dyn}(c) > \varepsilon_\mathrm{tol}$ is a
\emph{live holonomy violation}.
\end{definition}

\begin{definition}[Dynamic Order Parameter]
\label{def:dynamic-order}
The \emph{dynamic order parameter} at time $t$ is:
\[
  R_\mathrm{dyn}(t) = \abs[\bigg]{\frac{1}{n} \sum_{k=1}^n
  e^{i\hat{\theta}_k(t)}},
\]
where $\hat{\theta}_k(t)$ is $u_k$'s phase estimated from $\gamma_{u_k}$
at time $t$ (e.g., via Fourier analysis of the output time series).
\end{definition}

\begin{proposition}[Aperture Drift Detection]
\label{prop:aperture-drift}
A unit $u$ exhibits \emph{aperture drift} if $\gamma_u(t) \notin P_u$
for some $t \in [0,T]$ (the runtime trajectory exits the statically
declared production set). Aperture drift is a live specification violation.
\end{proposition}

\subsection{Purposelessness Detection}

\begin{definition}[Contribution Score]
\label{def:contribution}
The \emph{contribution score} of unit $u$ in ensemble $E$ is:
\[
  \delta\mathcal{S}(u, E) = \Sflat(E) - \Sflat(E \setminus \{u\}).
\]
$\delta\mathcal{S}(u, E) = 0$ iff $u$ is purposeless;
$\delta\mathcal{S}(u, E) > 0$ iff $u$ is purposeful (reduces the floor).
\end{definition}

Since $\Sflat(E)$ cannot be computed exactly (No Template Theorem), the
protocol approximates contribution scores via the \emph{OR-success
probability}:

\begin{definition}[OR-Success Probability]
\label{def:or-success}
Let $\mathcal{S}_\mathrm{thresh}$ be an entropy threshold. The
\emph{OR-success probability} of $E$ under load $\Lambda$ is:
\[
  P_\OR(E, \Lambda) = \Pr_{x \sim \Lambda}\!\bigl[\mathcal{S}(x)
  \leq \mathcal{S}_\mathrm{thresh}\bigr].
\]
Unit $u$ is \emph{practically purposeless} if
$P_\OR(E, \Lambda) \approx P_\OR(E \setminus \{u\}, \Lambda)$
up to significance level $\alpha$.
\end{definition}

\begin{definition}[Wind Tunnel Metric]
\label{def:wt-metric}
The \emph{wind tunnel metric} for ensemble $E$ under load $\Lambda$ is:
\[
  \mathrm{WT}(E, \Lambda) = \bigl(R_\mathrm{dyn},\,
  \mathcal{S}_\flat^\mathrm{est},\, \mathcal{H},\, \mathcal{D},\,
  \delta\mathcal{S}\bigr),
\]
where $R_\mathrm{dyn}$ is the dynamic order parameter time series,
$\mathcal{S}_\flat^\mathrm{est}$ is the estimated floor, $\mathcal{H}$
is the set of live holonomy violations, $\mathcal{D}$ is the set of
decoherence zones, and $\delta\mathcal{S}$ is the per-unit contribution
score vector.
\end{definition}

\begin{remark}
The wind tunnel metric is not binary. It is a \emph{regime map}: a
multi-dimensional characterisation of the ensemble's coordination state.
A system with $R_\mathrm{dyn} = 0.92$ (Coherent) and two identified
decoherence zones provides richer information than a ``97 of 100 tests
pass'' verdict: it locates the coordination deficits and quantifies the
distance to phase-locked operation.
\end{remark}

\begin{definition}[Wind Tunnel Protocol]
\label{def:wt-protocol}
Given system $S = (G, \{(A_u, P_u, \mathrm{val}_u)\})$ and load $\Lambda$,
the protocol proceeds:
\begin{enumerate}[leftmargin=*]
  \item \textit{Static analysis.} Extract $\{(A_u, P_u,
        \omega_u^\mathrm{est})\}$. Compute $\vartheta(u_i, u_j)$ for
        all edges. Compute $R_\mathrm{est}$; classify regime. Identify
        decoherence zones. Enumerate directed cycles; compute static
        residuals. Estimate $K_c$ and compare to structural coupling.
  \item \textit{Dynamic analysis.} Run $S$ under $\Lambda$; record
        $\{(\gamma_u, \Gamma_u, M_u)\}$. Compute $R_\mathrm{dyn}(t)$.
        Verify aperture compliance. Measure $\hol_\mathrm{dyn}(c)$ for
        all cycles. Compare $R_\mathrm{dyn}$ to static prediction.
  \item \textit{Purposelessness analysis.} Compute $P_\OR(E, \Lambda)$.
        For each $u$: ablate or resimulate $E \setminus \{u\}$; compute
        $P_\OR(E \setminus \{u\}, \Lambda)$. Report $\delta\mathcal{S}
        (u, E)$ for each unit.
  \item \textit{Output.} Return $\mathrm{WT}(E, \Lambda)$.
\end{enumerate}
\end{definition}

%===========================================================================
\section{Related Work}
\label{sec:related}
%===========================================================================

\subsection{Unit Testing}

Myers \cite{myers1979} established the vocabulary of unit testing: test
cases, oracles, coverage criteria. The oracle assumption --- that correct
output is known in advance --- is the bounded-observer assumption
formalised in the No Template Theorem. Unit testing is a degenerate case
of wind tunnel testing with ensemble size $n = 1$, where $R_\mathrm{ens}$
is trivially $1$ and all global coordination properties are absent. The
Local Invisibility Theorem (\cref{thm:local-invisibility}) and Isolation
Blindness (\cref{thm:isolation-blind}) show that unit testing cannot
detect failures outside the singleton regime.

\subsection{Property-Based Testing}

Claessen and Hughes \cite{claessen2000} introduced QuickCheck, which
generates test inputs from type specifications and checks user-defined
invariants. This generalises the test-case lattice but does not relax the
bounded-observer constraint: the invariant specification is itself a
bounded observer, subject to the No Template Theorem. Wind tunnel testing
subsumes property-based testing by treating invariants as aperture and
production set constraints, verified both statically and dynamically.

\subsection{Model Checking and Formal Verification}

Clarke and Emerson \cite{clarke1981} introduced CTL model checking, which
exhaustively verifies finite-state models against temporal specifications.
This achieves completeness within the model but faces state explosion and,
more fundamentally, the specification is a bounded observer subject to the
No Template Theorem: it cannot encode correct behaviour for systems of
greater complexity. Wind tunnel testing operates precisely in the regimes
where state explosion prevents exhaustive model checking.

\subsection{Chaos Engineering}

Basiri et al.\ \cite{basiri2016} formalised chaos engineering as
deliberate failure injection into live systems to observe resilience.
Chaos engineering is dynamic and ensemble-level --- closer in spirit to
wind tunnel testing than unit testing. The key gap is the absence of a
formal correctness criterion. The wind tunnel framework provides it: a
chaos experiment maintaining $R_\mathrm{dyn} \geq 0.95$ throughout
perturbation is formally resilient under the Phase-locked regime
definition.

\subsection{Distributed Systems Theory}

Fischer, Lynch, and Paterson \cite{flp1985} proved the impossibility of
consensus in asynchronous distributed systems with one faulty process.
Lamport \cite{lamport1982} established limits of agreement under
adversarial conditions. Both are no-go theorems in the spirit of the
No Template Theorem: structural limits on what distributed or observing
systems can determine. Wind tunnel testing operates within these limits
by not requiring consensus on a single correct output, but instead
measuring convergence toward the positive-volume region $C^*$.

\subsection{Collective Motion}

Vicsek et al.\ \cite{vicsek1995} showed that local alignment rules in
biological systems produce global order above a critical density. The
Vicsek and Kuramoto models belong to the same universality class for
large $n$ \cite{strogatz2000}, sharing the same discontinuous transition
structure at the critical coupling. The wind tunnel framework imports
this universality directly: the software synchronisation transition of
\cref{thm:partition-extinction} is in the same universality class.

%===========================================================================
\section{Discussion}
\label{sec:discussion}
%===========================================================================

\subsection{Implications for Testing Practice}

The wind tunnel framework implies a reorientation of testing practice.
Binary pass/fail verdicts should be supplemented with regime
classification: reporting $R_\mathrm{dyn}$ and decoherence zone severity
gives engineers a qualitatively richer picture. A system in the Turbulent
regime with 100\% unit-test coverage is less tested, in the semantically
relevant sense, than a system in the Coherent regime with 60\% coverage.

The contribution score $\delta\mathcal{S}(u, E)$ provides a novel
criterion for \emph{semantic dead code detection}: code that is executed
but contributes nothing to the ensemble's floor reduction. This is
orthogonal to syntactic dead code (never executed): a unit can be on the
hot path yet purposeless, or rarely invoked yet critical.

\subsection{Computational Complexity}

Computing $\mathrm{WT}(E, \Lambda)$ exactly requires $n+1$ ensemble
evaluations (the full ensemble and $n$ ablated variants). In practice,
approximations are available: random ablation, importance sampling for
$P_\OR$, and spectral approximations to $R_\mathrm{dyn}$. Directed cycle
enumeration is worst-case exponential in $\abs{V}$ but tractable for
sparse system graphs with $\abs{V} \leq 10^3$.

\subsection{Relationship to Statistical Mechanics}

Semantic entropy $\mathcal{S}$ is deliberately named by analogy with
Boltzmann's statistical entropy $S_B = k_B \ln \Omega$ \cite{boltzmann1877},
where $\Omega$ counts microstates consistent with the observed macrostate.
High $S_B$ means many microstates are consistent --- low knowledge of the
system. High $\mathcal{S}$ means many system states are consistent with
observed behaviour --- low confidence in purposeful operation. The analogy
is not rhetorical: the Onsager-Machlup variational principle
\cite{onsager1953} implies systems evolve to minimise free energy, which
includes an entropic term. A testing framework based on $\mathcal{S}$
inherits this variational structure: the system's most likely evolution
is toward the minimum of $\mathcal{S}$, i.e., toward $C^*$.

\subsection{Limitations}

First, the natural frequency $\omega_u$ is generally not computable from
source code and must be estimated from profiling data; estimation error
propagates into $K_c$ and regime classification.

Second, the action-cell $C^*$ must be specified by the system designer
via the purpose functional $\Phi_E$. The framework provides a formal
language for this specification but does not generate it automatically.
This reflects the No Template Theorem: purpose cannot be derived purely
from structure; it requires a semantic input external to the system.
Similarly, specifying $T_c^{\mathrm{spec}}$ for each directed cycle is
part of specifying the system's purpose functional; holonomy-based
consistency checking presupposes this specification.

Third, dynamic holonomy measurement requires closing loops in the
dependency graph. For cycles with traversal times of hours or days,
full holonomy measurement may require simulation rather than live
observation.

Fourth, the Kuramoto mean-field analysis is exact only in the
thermodynamic limit $n \to \infty$. For small ensembles ($n < 20$),
finite-size effects are significant and regime boundaries should be
treated as approximate.

\subsection{Open Problems}

\begin{enumerate}[leftmargin=*]
  \item \textit{Optimal coupling topology:} Given units with known
        $\{\omega_i\}$, what graph $G$ minimises $K_c$ for a target
        $R_\mathrm{ens}$?
  \item \textit{Purpose-cell inference:} Can $C^*$ be inferred from a
        finite trajectory sample without explicit specification? This
        requires learning the manifold structure of $\mathcal{S}$.
  \item \textit{Dynamic complexity:} What is the complexity of computing
        $R_\mathrm{dyn}$ to within $\varepsilon$ accuracy as a function
        of $n$, $\abs{E}$, and $T$?
  \item \textit{Multi-scale holonomy:} Extend the holonomy framework from
        node-level cycles to sub-ensemble-level cycles, providing a
        hierarchical consistency criterion.
  \item \textit{Specification inference:} Can cycle specifications
        $T_c^{\mathrm{spec}}$ be inferred or approximated from observed
        system behaviour, reducing the specification burden on the designer?
\end{enumerate}

%===========================================================================
\section{Conclusion}
\label{sec:conclusion}
%===========================================================================

We have presented the wind tunnel framework for software correctness,
grounded in three foundational theorems.

The \emph{Local Invisibility Theorem} (\cref{thm:local-invisibility})
shows that global semantic failure is structurally undetectable by local
queries: a system can pass every unit test while accumulating global
inconsistency in its constraint cycles.

The \emph{No Template Theorem} (\cref{thm:no-template}) shows that no
bounded observer can encode a complete correctness criterion for a system
of equal or greater complexity. This makes template-matching an
epistemologically untenable foundation for correctness beyond a complexity
threshold.

The \emph{Cycle Inconsistency Theorem} (\cref{thm:hac}) provides a
principled detection criterion: nonzero holonomy on any directed cycle ---
measured as deviation from the cycle's intended transformation
$T_c^{\mathrm{spec}}$ --- implies global semantic incorrectness. Zero
holonomy is necessary but not sufficient; DAG-structured systems are
trivially holonomy-free, and their global correctness must be established
through boundary conditions rather than cycle consistency.

Building on these foundations, we defined semantic entropy $\mathcal{S}$
and the action-cell $C^*$ as rigorous substitutes for the binary oracle,
and the ensemble order parameter $R_\mathrm{ens}$ as a measure of global
coordination grounded in Kuramoto theory. We classified five coordination
regimes, proved the discontinuous nature of the synchronisation transition
(partition extinction), and established by \cref{thm:isolation-blind}
that purposelessness cannot be detected by unit isolation.

The resulting wind tunnel protocol is a two-phase instrument: a static
phase classifying the system's coordination potential from its structure,
and a dynamic phase measuring actual coherence under load. The primary
output is not a pass/fail verdict but a regime map --- a quantitative
characterisation of coordination health across the dependency graph.

The framework does not abolish unit testing; it contextualises it. Unit
tests are local queries that usefully catch aperture mismatches and
production set violations. But they are blind, by theorem, to the global
semantic properties that determine whether the assembled system fulfils
its purpose. The wind tunnel sees what the unit test structurally cannot:
the whole assembled system under real air.

\bigskip

\noindent\textbf{Acknowledgements.} The authors thank the mathematical
physics and distributed systems communities whose foundational results ---
particularly those of Kuramoto, Strogatz, Kirchhoff, Banach, and Shannon
--- made this synthesis possible.

\bibliographystyle{plainnat}
\bibliography{references}

\end{document}
